% DERIVED from formal_layer.json — read-only, do not edit.
\begin{assumptionv}{ass:iid-sampling}[Independent and identically distributed sampling]
The observed records \(O_1,\ldots,O_n\) are independent and identically distributed according to the observed-data law induced by \(P\), the full-data probability law.
\end{assumptionv}

\begin{definitionv}{synth_2}[Full-data record coordinates]
A full-data record, for alphabet size \(d\in\mathbb N\), consists of a baseline category \(X\), a binary arm indicator \(A\), control and treated surrogate coordinates \(S(0)\) and \(S(1)\), control and treated outcome coordinates \(Y(0)\) and \(Y(1)\), and an arrival flag \(R\), with domains
\[
X\in\{x\in\mathbb Z:0\le x<d\},
\qquad
\bigl(A,S(0),S(1),Y(0),Y(1),R\bigr)\in\{0,1\}^{6}.
\]
\end{definitionv}

\begin{definitionv}{synth_3}[Potential and assigned surrogates]
The potential surrogate \(S(a)\), for \(a\in\{0,1\}\), denotes the record's binary control surrogate when \(a=0\) and binary treated surrogate when \(a=1\). The assigned-arm surrogate \(S\), for any full-data record with binary assignment \(A\in\{0,1\}\), is defined by
\[
S=
\begin{cases}
S(1), & A=1,\\
S(0), & A=0.
\end{cases}
\]
\end{definitionv}

\begin{definitionv}{synth_4}[Potential and assigned-arm outcomes]
The binary potential outcomes \(Y(0)\) and \(Y(1)\) are the control and treated outcome coordinates of a full-data record. For a nonnegative integer \(d\), such a record consists of a baseline category \(X\) in an alphabet of size \(d\) and binary coordinates \(A,S(0),S(1),Y(0),Y(1),R\). The assigned-arm outcome \(Y\) is defined by
\[
Y=
\begin{cases}
Y(1), & A=1,\\
Y(0), & A=0.
\end{cases}
\]
The space of full-data records is equipped with the measurable structure in which every subset is measurable.
\end{definitionv}

\begin{assumptionv}{ass:randomized-independence}[Independence of treatment assignment]
Under \(P\), the full-data probability law, treatment assignment \(A\) satisfies
\[
A\perp\!\!\!\perp\{X,S(0),S(1),Y(0),Y(1)\}.
\]
Here \(X\) is the baseline covariate, \(S(0),S(1)\) are the potential surrogates, and \(Y(0),Y(1)\) are the potential outcomes.
\end{assumptionv}

\begin{assumptionv}{ass:balanced-randomization}[Balanced treatment assignment]
Under \(P\), the full-data probability law, treatment assignment \(A\) satisfies
\[
P(A=1)=\frac12.
\]
\end{assumptionv}

\begin{definitionv}{synth_1}[Full-data probability law]
The full-data law \(P\), for a nonnegative integer \(d\), is a probability measure on the finite record space
\[
(X,A,S(0),S(1),Y(0),Y(1),R)
\in \{0,\ldots,d-1\}\times\{0,1\}^{6}.
\]
Here \(X\) is the baseline category, \(A\) is the treatment arm, \(S(0)\) and \(S(1)\) are the potential surrogates, \(Y(0)\) and \(Y(1)\) are the potential outcomes, and \(R\) is the arrival flag. The baseline alphabet is empty when \(d=0\).
\end{definitionv}

\begin{assumptionv}{ass:surrogate-consistency}[Consistency of the surrogate]
The realized surrogate \(S\) equals the potential surrogate corresponding to treatment assignment \(A\):
\[
S=S(A)\qquad\text{almost surely under }P,
\]
where \(P\) is the full-data probability law.
\end{assumptionv}

\begin{assumptionv}{ass:outcome-consistency}[Consistency of the outcome]
The realized outcome \(Y\) equals the potential outcome corresponding to treatment assignment \(A\):
\[
Y=Y(A)\qquad\text{almost surely under }P,
\]
where \(P\) is the full-data probability law.
\end{assumptionv}

\begin{assumptionv}{ass:arrival-mar}[Conditionally independent outcome arrival]
Under \(P\), the full-data probability law, the outcome-arrival indicator \(R\) and realized outcome \(Y\) are conditionally independent given baseline \(X\), treatment assignment \(A\), and realized surrogate \(S\):
\[
R\perp\!\!\!\perp Y\mid(X,A,S).
\]
\end{assumptionv}

\begin{definitionv}{synth_6}[Cell alphabet and membership]
The cell alphabet \(\mathcal J_d\), for a nonnegative integer \(d\), is
\[
\mathcal J_d
=
\{0,1\}\times\{x\in\mathbb Z:0\le x<d\}\times\{0,1\}.
\]
For a full-data record with baseline category \(X\), binary assignment \(A\), and realized surrogate \(S\), membership in a cell \(j=(a,x,s)\in\mathcal J_d\) is defined by
\[
A=a \quad\land\quad X=x \quad\land\quad S=s.
\]
\end{definitionv}

\begin{definitionv}{synth_16}[Cell probabilities \(p_{axs}(P)\) and \(\rho_{axs}(P)\)]
The full-membership probability of a cell is
\[
p_{axs}(P):=P(A=a,X=x,S=s),
\qquad
(a,x,s)\in\mathcal J_d
=\{0,1\}\times\{0,\ldots,d-1\}\times\{0,1\},
\]
where \(P\) is a full-data law. On every occupied cell, meaning \(p_{axs}(P)>0\), define its conditional outcome-arrival probability by
\[
\rho_{axs}(P)
:=
P(R=1\mid A=a,X=x,S=s)
=
\frac{P(A=a,X=x,S=s,R=1)}{p_{axs}(P)}.
\]
\end{definitionv}

\begin{assumptionv}{ass:occupied-cell-arrival}[Occupied-cell arrival floor]
For each cell
\[
(a,x,s)\in\mathcal J_d
=\{0,1\}\times\{0,\ldots,d-1\}\times\{0,1\},
\]
let \(p_{axs}(P)\), the full-membership cell probability, and \(\rho_{axs}(P)\), the conditional outcome-arrival probability on an occupied cell, be defined as in \cref{obj:synth_16}. The occupied-cell arrival condition is
\[
p_{axs}(P)>0
\quad\Longrightarrow\quad
\rho_{axs}(P)\ge q.
\]
\end{assumptionv}

\begin{definitionv}{def:unrestricted-arrival-model-class}[Arrival-floor model \(\mathcal M^{+}_{n,d,q}\)]
The class \(\mathcal M^{+}_{n,d,q}\) consists of full-data laws \(P\), for integers \(n,d\ge1\) and an arrival floor \(0<q\le1\), satisfying the following conditions:
\begin{itemize}
\item \textbf{(Randomized independence.)} \cref{obj:ass:randomized-independence}.
\item \textbf{(Balanced randomization.)} \cref{obj:ass:balanced-randomization}.
\item \textbf{(Surrogate consistency.)} \cref{obj:ass:surrogate-consistency}.
\item \textbf{(Outcome consistency.)} \cref{obj:ass:outcome-consistency}.
\item \textbf{(Arrival MAR.)} \cref{obj:ass:arrival-mar}.
\item \textbf{(Occupied-cell arrival.)} \cref{obj:ass:occupied-cell-arrival}.
\end{itemize}
The arrival floor ranges over the entire interval \(0<q\le1\). Independent, identically distributed sampling is imposed separately on the experiment through \cref{obj:ass:iid-sampling}.
\end{definitionv}

\begin{definitionv}{def:ate-functional}[Average treatment effect \(\tau(P)\)]
The average treatment effect is
\[
\tau(P):=E_P[Y(1)-Y(0)],
\]
where \(Y(1)\) and \(Y(0)\) are the potential outcomes under treatment and control.
\end{definitionv}

\begin{definitionv}{synth_19}[Treatment sign function]
Define the treatment sign function \(g:\{0,1\}\to\{-1,1\}\) by
\[
g(a):=2a-1,\qquad a\in\{0,1\}.
\]
Thus \(g(0)=-1\) and \(g(1)=1\).
\end{definitionv}

\begin{definitionv}{def:cell-functional}[Signed cell functional \(\Phi(P)\)]
The total signed observed-cell functional is
\[
\Phi(P):=2\sum_{(a,x,s)\in\mathcal J_d}g(a)c_{axs}(P),
\]
where \(\mathcal J_d\) is the treatment–baseline–surrogate cell-index set and \(g(a)\) is the treatment sign function. On this index set, with \(x\in\{0,\ldots,d-1\}\), the full-membership cell probability \(p_{axs}(P)\) is defined in \cref{obj:synth_16}.

For a cell with positive arrived-cell probability,
\[
P(A=a,X=x,S=s,R=1)>0,
\]
define its arrived-cell outcome mean and weighted contribution by
\[
\mu_{axs}(P):=E_P[Y\mid A=a,X=x,S=s,R=1],
\qquad
c_{axs}(P):=p_{axs}(P)\mu_{axs}(P).
\]
When the arrived-cell probability is zero, including whenever \(p_{axs}(P)=0\), set
\[
\mu_{axs}(P):=0,
\qquad
c_{axs}(P):=0.
\]
Thus \(\Phi\) is defined on every full-data law.
\end{definitionv}

\begin{definitionv}{synth_5}[Observed sample and its canonical law]
The random observed sample is \(\mathbf O_n=(O_1,\ldots,O_n)\), where \(O_i=(X_i,A_i,S_i,R_i,R_iY_i)\). For nonnegative integers \(n,d\) and a full-data probability law \(P\) with \(d\) baseline categories, its canonical law is
\[
\mathcal L_P(\mathbf O_n):=\bigl(\mathcal L_P(X,A,S,R,RY)\bigr)^{\otimes n}.
\]
Here \(\mathcal L_P(X,A,S,R,RY)\) is the distribution under \(P\) of the observed record \((X,A,S,R,RY)\), whose last four coordinates are binary and whose final coordinate is the product of arrival and outcome. Thus \(\mathcal L_P(\mathbf O_n)\) is the law of \(n\) independent observed records with this common distribution.
\end{definitionv}

\begin{definitionv}{def:unrestricted-minimax-risk}[Unrestricted minimax squared-error risk]
The minimax squared-error risk \(\mathfrak R^{+}_{n,d,q}\), for \(n,d\in\mathbb N\) and \(q\in\mathbb R\), is defined by
\[
\mathfrak R^{+}_{n,d,q}
:=
\inf_{\mathsf T\in\mathcal T_n}
\sup_{P\in\mathcal M^{+}_{n,d,q}}
\int\!\int_{[-1,1]}
(h-\tau(P))^2\,
\mathsf T(o,dh)\,
\mathcal L_P(\mathbf O_n)(do).
\]
Here \(\mathcal M^{+}_{n,d,q}\) is the full-data model class in \cref{obj:def:unrestricted-arrival-model-class}, and \(\tau(P)\) is the population average treatment effect in \cref{obj:def:ate-functional}. The class \(\mathcal T_n\) consists of all parameter-independent Markov kernels from the observed sample space to \([-1,1]\), including deterministic estimators as Dirac kernels.

The random observed-record tuple \(\mathbf O_n\) and its product law under the sampling in \cref{obj:ass:iid-sampling} are
\[
\mathbf O_n=(O_1,\ldots,O_n),
\qquad
\mathcal L_P(\mathbf O_n)
=
\bigl(\mathcal L_P(O_1)\bigr)^{\otimes n},
\]
where each \(O_i\) is an observed record induced by \(P\). The risk integrates over both the observed sample and the kernel draw.

By \cref{obj:lem:randomized-kernel-barycenter}, the defining infimum equals the infimum over deterministic estimators taking values in \([-1,1]\). At \(q=1\), the same decision class \(\mathcal T_n\) defines the minimax risk over \(\mathcal M^{+}_{n,d,1}\), whose laws satisfy \(R=1\) almost surely.
\end{definitionv}

\begin{definitionv}{synth_7}[Total real logarithmic scale]
The logarithmic scale \(\ell\), for \(n\in\mathbb N\) and \(q\in\mathbb R\), is defined by
\[
\ell=
\begin{cases}
\log\bigl|\exp(1)+nq\bigr|, & \exp(1)+nq\ne 0,\\
0, & \exp(1)+nq=0,
\end{cases}
\]
where \(\log\) on positive arguments is the natural logarithm.
\end{definitionv}

\begin{assumptionv}{ass:rare-arrival-slice}[Rare-arrival restriction]
The arrival floor \(q\) and logarithmic scale \(\ell\) satisfy
\[
q\ell^2\le \frac{1}{64}.
\]
\end{assumptionv}

\begin{definitionv}{def:model-class}[Rare-arrival model \(\mathcal M_{n,d,q}\)]
\(\mathcal M_{n,d,q}\) is the class of full-data probability laws \(P\) satisfying the following conditions:
\begin{itemize}
\item \textbf{(Parameters.)} \(n,d\ge1\) and \(0<q\le1\).
\item \textbf{(Sampling.)} \cref{obj:ass:iid-sampling} holds.
\item \textbf{(Randomization.)} \cref{obj:ass:randomized-independence,obj:ass:balanced-randomization} hold.
\item \textbf{(Consistency.)} \cref{obj:ass:surrogate-consistency,obj:ass:outcome-consistency} hold.
\item \textbf{(Arrival.)} \cref{obj:ass:arrival-mar,obj:ass:occupied-cell-arrival} hold.
\item \textbf{(Rare-arrival restriction.)} \cref{obj:ass:rare-arrival-slice} holds.
\end{itemize}
\end{definitionv}

\begin{definitionv}{def:observed-experiment}[Observed sample experiment]
The observed experiment \(\mathcal E_{n,d,q}\), for \(n,d\in\mathbb N\) and \(q\in\mathbb R\), is the family
\[
\mathcal E_{n,d,q}
:=
\bigl\{\mathcal L_P(\mathbf O_n):P\in\mathcal M_{n,d,q}\bigr\},
\]
where \(\mathcal M_{n,d,q}\) is the model class defined in \cref{obj:def:model-class}. Here \(\mathbf O_n\), the random observed-record tuple, and its law are specified by
\[
\mathbf O_n=(O_i)_{i=1}^{n},
\qquad
O_i=(X_i,A_i,S_i,R_i,R_iY_i),
\qquad
\mathcal L_P(\mathbf O_n)
=
\bigl(\mathcal L_P(X,A,S,R,RY)\bigr)^{\otimes n}.
\]
Thus the records are independent and identically distributed with the observed marginal law induced by \(P\). Each record retains the baseline category \(X\in\{0,\ldots,d-1\}\), assignment \(A\), realized surrogate \(S\), arrival indicator \(R\), and arrival-multiplied outcome \(RY\); the last four coordinates are binary. The record space carries the discrete sigma-algebra.
\end{definitionv}

\begin{definitionv}{def:minimax-risk}[Minimax squared-error risk]
The minimax squared-error risk \(\mathfrak R_{n,d,q}\), for \(n,d\in\mathbb N\) and \(q\in\mathbb R\), is
\[
\mathfrak R_{n,d,q}
:=
\inf_{\mathsf T\in\mathcal T_n}
\sup_{P\in\mathcal M_{n,d,q}}
\int\!\int_{[-1,1]}
(h-\tau(P))^2\,
\mathsf T(o,dh)\,
\mathcal L_P(\mathbf O_n)(do).
\]
Here \(\mathcal M_{n,d,q}\) is the full-data model class of \cref{obj:def:model-class}, and \(\tau(P)\) is the population average treatment effect of \cref{obj:def:ate-functional}. The observed sample \(\mathbf O_n=(O_i)_{i=1}^n\) denotes the random tuple of observed records. Its law, under \cref{obj:ass:iid-sampling}, is
\[
\mathcal L_P(\mathbf O_n)
=
\bigotimes_{i=1}^{n}\mathcal L_P(O_i).
\]
The decision class \(\mathcal T_n\) consists of all parameter-independent Markov kernels from the observed-sample space to \([-1,1]\), including deterministic estimators represented by Dirac kernels.

The notation \(E_P[(\mathsf T-\tau(P))^2]\) denotes the displayed joint sample-and-kernel risk. By \cref{obj:lem:randomized-kernel-barycenter}, the minimax value also equals the infimum of the maximal squared-error risk over deterministic estimators taking values in \([-1,1]\).
\end{definitionv}

\begin{definitionv}{def:frontier-rate}[Frontier rate \(r_{n,d,q}\)]
The frontier rate is
\[
r_{n,d,q}
:=
\min\left\{
1,\,
\frac{1}{nq}
+
\left[\frac{d}{nq\log(e+nq)}\right]^2
\right\}.
\]
\end{definitionv}

\begin{definitionv}{def:complete-arrival-ht-estimator}[Clipped observed-outcome arm contrast]
The clipped observed-outcome arm contrast is the map
\[
\widehat\tau^{\mathrm{HT}}_n(\mathbf O_n)
:=
\Pi_{[-1,1]}\!\left\{
\frac{2}{n}\sum_{i=1}^n
(2A_i-1)\mathbf{1}\{R_i=1,\ R_iY_i=1\}
\right\},
\qquad n\ge1.
\]
Here \(n,d\) are nonnegative integers, and \(\mathbf O_n\), the observed sample, consists of records
\(O_i=(X_i,A_i,S_i,R_i,R_iY_i)\), with \(X_i\in\{0,\ldots,d-1\}\) and all four remaining coordinates binary; the baseline support is empty when \(d=0\). The notation \(R_iY_i\) denotes the stored arrived-outcome coordinate; the indicator equals one exactly when both this coordinate and the arrival coordinate \(R_i\) equal one. Projection and clipping coincide:
\[
\Pi_{[-1,1]}(u)=\operatorname{clip}_{[-1,1]}(u)
=\max\{-1,\min\{1,u\}\},
\qquad u\in\mathbb R.
\]
For \(n=0\), the map is defined to equal zero.
\end{definitionv}

\begin{definitionv}{synth_8}[Polynomial degree tuning]
The degree \(k\), for \(n\in\mathbb N\) and \(q\in\mathbb R\), is defined by
\[
k=
\begin{cases}
\displaystyle \max\left\{0,\left\lfloor\frac{\log|e+nq|}{64}\right\rfloor\right\},
& e+nq\ne 0,\\
0, & e+nq=0.
\end{cases}
\]
\end{definitionv}

\begin{definitionv}{synth_18}[Chebyshev needle polynomial]
Put \(N=nq\), \(\ell=\log(e+N)\), \(k=\lfloor\ell/64\rfloor\), and \(B=256\ell\). Let \(T_k\) be the first-kind Chebyshev polynomial, characterized by \(T_k(\cos u)=\cos(ku)\). On the branch \(\ell\ge128\) and \(d\ge\sqrt N\), where \(k\ge2\), define
\[
Q_k(t):=
\begin{cases}
\displaystyle\frac{1-T_k(1-2t/B)}{2k^2t/B},&t>0,\\[6pt]
1,&t=0.
\end{cases}
\]
Here \(t\ge0\) is an arrived-count intensity, and the value at zero is the continuous polynomial extension of the displayed quotient. On every other branch, set \(Q_k(t):=1\) for every \(t\ge0\).
\end{definitionv}

\begin{definitionv}{synth_13}[Polynomial coefficient rule]
The coefficients \(a_v\), for \(n,d,v\in\mathbb N\) and \(q\in\mathbb R\), are defined by
\[
a_v=
\begin{cases}
[t^v]\bigl(1-Q_k(t)\bigr),
& 128\le \ell \ \text{and}\ \sqrt{N}\le d,\\
0,
& \text{otherwise}.
\end{cases}
\]
Here \([t^v]\) denotes the coefficient of \(t^v\), \(N=nq\) is the effective sample size, \(\ell\) is the logarithmic scale, and \(Q_k\) is the Chebyshev needle polynomial. On the indicated branch, this polynomial is the needle radius divided by \(2k^2\), multiplied by the polynomial quotient of
\[
1-T_k\!\left(1-\frac{2t}{\text{needle radius}}\right)
\]
upon division by \(t\), with the remainder discarded. The degree \(k\) and needle radius are the construction's values at \((n,q)\), and \(T_k\) is the first-kind Chebyshev polynomial.
\end{definitionv}

\begin{definitionv}{synth_17}[Arrived-outcome counts \(C'_j\) and \(C_j\)]
The pilot and estimation arrived-outcome counts are
\[
C'_j:=\sum_{i\in\mathcal I_{\mathrm{pil}}}
\mathbf1\{A_i=a,\ X_i=x,\ S_i=s,\ R_i=1\},
\qquad
C_j:=\sum_{i\in\mathcal I_{\mathrm{est}}}
\mathbf1\{A_i=a,\ X_i=x,\ S_i=s,\ R_i=1\},
\quad j=(a,x,s)\in\mathcal J_d.
\]
Here \(\mathbf O_n\), the random observed sample, consists of the records
\[
O_i=(X_i,A_i,S_i,R_i,R_iY_i),
\qquad
\mathbf O_n=(O_i)_{i=1}^n,
\qquad
\mathcal L_P(\mathbf O_n)=\bigl(\mathcal L_P(O_1)\bigr)^{\otimes n}.
\]
Given a prefix length \(K_0\le n\) and an assignment of its indices to membership, pilot, and estimation streams, \(\mathcal I_{\mathrm{pil}}\) and \(\mathcal I_{\mathrm{est}}\) denote the sets of prefix indices assigned to the pilot and estimation streams, respectively. The cell alphabet and observed-sample law are defined in \cref{obj:synth_6,obj:synth_5}; the prefix and stream assignments are those used in \cref{obj:def:mixed-count-estimator}.
\end{definitionv}

\begin{definitionv}{synth_10}[Estimation arrived-one count]
The estimation arrived-one count \(U_j\) is defined for any \(k,d\in\mathbb N\), collection of \(k\) observed records, assignment of each record to a stream in \(\{0,1,2\}\), and cell \(j=(a,x,s)\in\{0,1\}\times\{0,\ldots,d-1\}\times\{0,1\}\) by
\[
U_j
:=
\#\bigl\{\,0\le i<k:
i\text{ is assigned to stream }2,\ 
A_i=a,\ X_i=x,\ S_i=s,\ R_i=1,\ (RY)_i=1
\,\bigr\}.
\]
Here \(O_i=(X_i,A_i,S_i,R_i,(RY)_i)\) denotes the \(i\)th observed record, with baseline coordinate \(X_i\in\{0,\ldots,d-1\}\) and binary coordinates \(A_i,S_i,R_i,(RY)_i\).
\end{definitionv}

\begin{definitionv}{synth_15}[Selected polynomial cell statistic]
The selected cell statistic \(G_j\), for natural-number parameters \(n,d\), a real parameter \(q\), a finite collection of observed records with baseline coordinate in \(\{0,\ldots,d-1\}\) and binary assignment, surrogate, arrival, and arrived-outcome coordinates, an assignment of each record to one of three streams, and a cell \(j\in\{0,1\}\times\{0,\ldots,d-1\}\times\{0,1\}\), is defined by
\[
G_j=
\begin{cases}
\displaystyle
\sum_{v=1}^{\max\{k-1,0\}}
a_v U_j
\prod_{t=0}^{v-2}\max\{C_j-1-t,0\},
&
128\le \ell,\quad \sqrt{N}\le d,\quad C'_j\le B/4,
\\[6pt]
\displaystyle U_j/C_j,
&
C_j>0
\ \text{and the preceding branch condition fails},
\\[4pt]
0,
&
C_j=0
\ \text{and the preceding branch condition fails}.
\end{cases}
\]
Here \(N=nq\) is the effective sample size, \(\ell\) is the construction's logarithmic scale, \(B=256\ell\) is its threshold, \(k\) is its polynomial degree, and \(a_v\) are its polynomial coefficients. The quantities \(C'_j\), \(C_j\), and \(U_j\) are, respectively, the pilot-stream arrived count, the estimation-stream arrived count, and the estimation-stream arrived count with outcome equal to one in cell \(j\). An empty sum equals zero and an empty product equals one.
\end{definitionv}

\begin{definitionv}{synth_14}[Factorial cell statistic]
The cell statistic \(H_j\), for a finite collection of observed records, assignments to streams \(0,1,2\), a cell \(j\in\{0,1\}\times\{0,\ldots,d-1\}\times\{0,1\}\), \(n,d\in\mathbb N\), and \(q\in\mathbb R\), is
\[
H_j=
\begin{cases}
\displaystyle
\sum_{v=1}^{\max\{k-1,0\}}
a_v U_j
\prod_{r=1}^{v-1}\max\{C_j-r,0\},
& \ell\ge128\ \text{and}\ \sqrt N\le d,\\[6pt]
0,&\text{otherwise}.
\end{cases}
\]
Here \(N=nq\) is the effective sample size, \(\ell\) is the logarithmic scale, and the degree parameter is
\[
k=\max\{\lfloor\ell/64\rfloor,0\}.
\]
The estimation-stream arrival count \(C_j\) and arrived-one count \(U_j\) are
\[
C_j=\sum_i
\mathbf 1\{\text{record }i\text{ is assigned to stream }2,\ 
(A_i,X_i,S_i)=j,\ R_i=1\},
\qquad
U_j=\sum_i
\mathbf 1\{\text{record }i\text{ is assigned to stream }2,\ 
(A_i,X_i,S_i)=j,\ R_i=1,\ (RY)_i=1\}.
\]
The coefficients associated with the Chebyshev polynomial \(Q_k\) are
\[
a_v=
\begin{cases}
[t^v]\bigl(1-Q_k(t)\bigr),
&\ell\ge128\ \text{and}\ \sqrt N\le d,\\
0,&\text{otherwise},
\end{cases}
\]
where \([t^v]\) denotes the coefficient of \(t^v\). Empty sums equal zero and empty products equal one.
\end{definitionv}

\begin{definitionv}{synth_12}[Elementary cell ratio statistic]
The elementary cell statistic \(D_j\), for any finite collection of observed records, any assignment of those records to three streams indexed by \(0,1,2\), and any cell \(j\in\{0,1\}\times\{0,\ldots,d-1\}\times\{0,1\}\), is
\[
D_j=
\begin{cases}
U_j/C_j, & C_j>0,\\
0, & C_j=0.
\end{cases}
\]
Here the records have baseline coordinate \(X\in\{0,\ldots,d-1\}\) and binary coordinates \(A,S,R,RY\), and the counts are
\[
\begin{aligned}
C_j&=\#\{\text{records assigned to stream }2:
                 (A,X,S)=j,\ R=1\},\\
U_j&=\#\{\text{records assigned to stream }2:
                 (A,X,S)=j,\ R=1,\ RY=1\}.
\end{aligned}
\]
\end{definitionv}

\begin{definitionv}{synth_9}[Membership stream cell count]
The membership count \(M_j\), for any nonnegative integers \(k,d\), a collection of \(k\) observed records \(O_i\) with baseline coordinates in \(\{0,\ldots,d-1\}\) and binary assignment, surrogate, arrival, and recorded-outcome coordinates, an assignment of each record to one of the streams \(0,1,2\), and a cell \(j=(a,x,s)\in\{0,1\}\times\{0,\ldots,d-1\}\times\{0,1\}\), is
\[
M_j
:=
\#\bigl\{i:0\le i<k,\ 
i\text{ is assigned to stream }0,\ 
A_i=a,\ X_i=x,\ S_i=s\bigr\}.
\]
\end{definitionv}

\begin{definitionv}{synth_11}[Cell membership weight]
The cell-weight statistic \(W_j\) is defined for nonnegative integers \(k,d,n\), a collection of \(k\) observed records \(O_i\) with baseline coordinate \(X_i\in\{0,\ldots,d-1\}\) and binary coordinates \(A_i,S_i,R_i,RY_i\), an assignment of each record to a stream labelled \(0,1,2\), and a cell
\[
j\in\{0,1\}\times\{0,\ldots,d-1\}\times\{0,1\}.
\]
Using the membership count and stream intensity
\[
M_j=\#\{i\in\{0,\ldots,k-1\}: O_i\text{ is assigned to stream }0,\ (A_i,X_i,S_i)=j\},
\qquad
m=\frac{n}{6},
\]
define
\[
W_j=\frac{M_j}{m}.
\]
The count \(M_j\) is interpreted as a real number in this quotient, and division by zero is assigned the value zero, so \(W_j=0\) when \(n=0\). The index sets \(\{0,\ldots,k-1\}\) and \(\{0,\ldots,d-1\}\) are empty when their respective sizes are zero.
\end{definitionv}

\begin{algorithmv}{def:mixed-count-estimator}[Mixed-count estimator $\widehat\tau^{\mathrm{mix}}_{n,d,q}$]
Given the observed sample \(\mathbf O_n\), parameters \(n,d,q\), effective sample size \(N=nq\), logarithmic scale \(\ell\), and the declared cell-statistic and polynomial formulas, define the estimator by the following steps.
\begin{enumerate}
\item Draw an auxiliary prefix length independently of the data:
\[
K_0\sim\operatorname{Poisson}(n/2).
\]
When \(K_0\le n\), independently assign each of the first \(K_0\) observations to the membership, pilot, or estimation stream, each with probability \(1/3\). Denote the prefix and stream-assignment randomization by \(\omega\).

\item On the event \(K_0>n\), set
\[
\widetilde\tau(\mathbf O_n;\omega)=0.
\]
When \(K_0\le n\), form the declared statistics \(M_j,C'_j,C_j,U_j,W_j,D_j\). If \(N\le1\), set
\[
\widetilde\tau(\mathbf O_n;\omega)=0.
\]

\item When \(K_0\le n\), \(N>1\), and either \(\ell<128\) or \(d<\sqrt N\), set
\[
\widetilde\tau(\mathbf O_n;\omega)
=
\operatorname{clip}_{[-1,1]}
\left\{
2\sum_{j=(a,x,s)}g(a)W_jD_j
\right\}.
\]
On this branch, the Chebyshev quotient and its coefficients are left unevaluated.

\item When \(K_0\le n\), \(N>1\), \(\ell\ge128\), and \(d\ge\sqrt N\), form the declared total quantities \(Q_k,a_v,H_j,G_j\) and set
\[
\widetilde\tau(\mathbf O_n;\omega)
=
\operatorname{clip}_{[-1,1]}
\left\{
2\sum_{j=(a,x,s)}g(a)W_jG_j
\right\}.
\]

\item Output the conditional average
\[
\widehat\tau^{\mathrm{mix}}_{n,d,q}(\mathbf O_n)
:=
E_\omega\!\left[
\widetilde\tau(\mathbf O_n;\omega)\mid\mathbf O_n
\right].
\]
This average is a finite sum over prefix lengths \(K_0\le n\) and their stream assignments, together with the zero contribution from the Poisson tail. It defines a total, deterministic, sample-measurable estimator whose squared-error risk is no larger than that of the auxiliary randomized output.
\end{enumerate}
\end{algorithmv}

\begin{theoremv}{thm:unrestricted-near-complete-arrival-envelope}[Near-complete arrival risk envelope]
There exists a universal constant \(\overline C>0\) with the following properties. For every \(n,d\ge1\) and \(0<q\le1\), let \(r_{n,d,q}\) be the rate in \cref{obj:def:frontier-rate}, let \(\ell\) be the logarithmic scale in \cref{obj:synth_7}, let \(k\) be the polynomial degree specified in \cref{obj:synth_8}, and let \(B\) be the threshold specified in \cref{obj:synth_18}. Define the effective sample size \(N\), stream intensity \(m\), effective stream size \(N_0\), and prefix-tail constant \(\gamma_0\) by
\[
N=nq,\qquad m=\frac n6,\qquad N_0=mq,\qquad
\gamma_0=\log 2-\frac12.
\]
Define the deterministic squared-error envelope \(\mathfrak U_{n,d,q}\) as follows. When \(N\le1\), set
\[
\mathfrak U_{n,d,q}=1.
\]
When \(N>1\) and the estimator in \cref{obj:def:mixed-count-estimator} uses its elementary branch, set
\[
\mathfrak U_{n,d,q}
=
\min\left\{
4,\,
\left(\frac{8d}{eN_0}\right)^2
+4\left(\frac4{N_0}+\frac1m\right)
+4e^{-n\gamma_0}
\right\}.
\]
When \(N>1\) and that estimator uses its polynomial branch, define the certificate quantities
\[
b_{n,d,q}
=
2\left(\frac{4dB}{N_0k^2}+3e^{-16\ell}\right),
\]
\[
v_{n,d,q}
=
8\left(\frac4{N_0}+\frac1m\right)
+64d\left(e^{\ell/8}+1\right)
\left(\frac{B^2}{N_0^2}+\frac{B}{mN_0}\right)
+32e^{-32\ell}\left(1+\frac1m\right),
\]
and set
\[
\mathfrak U_{n,d,q}
=
\min\left\{
4,\,
b_{n,d,q}^{\,2}+v_{n,d,q}+4e^{-n\gamma_0}
\right\}.
\]

For every full-data law \(P\in\mathcal M^{+}_{n,d,q}\), the unrestricted-arrival model class in \cref{obj:def:unrestricted-arrival-model-class}, suppose that \(\mathbf O_n\), the observed sample, satisfies \cref{obj:ass:iid-sampling} under the canonical product law of \(n\) observed records induced by \(P\). The mixed-count estimator of \cref{obj:def:mixed-count-estimator} satisfies
\[
E_P\!\left[
\left\{
\widehat\tau^{\mathrm{mix}}_{n,d,q}(\mathbf O_n)-\tau(P)
\right\}^{2}
\right]
\le \mathfrak U_{n,d,q}
\le \overline C\,r_{n,d,q},
\]
where \(\tau(P)\) is the average treatment effect in \cref{obj:def:ate-functional}. It also satisfies
\[
E_P\!\left[
\left\{
\widehat\tau^{\mathrm{mix}}_{n,d,q}(\mathbf O_n)-\tau(P)
\right\}^{2}
\right]
\le
E_P\!\left[
E_\omega\!\left[
\left\{
\widetilde\tau(\mathbf O_n;\omega)-\tau(P)
\right\}^{2}
\right]
\right],
\]
where \(\widetilde\tau(\mathbf O_n;\omega)\) is the auxiliary randomized output in \cref{obj:def:mixed-count-estimator}, and the inner expectation uses its specified auxiliary randomization, including the zero output when the Poisson prefix length exceeds \(n\).

For every \(n,d\ge1\) and \(0<q\le1\), the minimax squared-error risk \(\mathfrak R^{+}_{n,d,q}\) over the parameter-independent bounded estimator kernels in \cref{obj:def:unrestricted-minimax-risk} obeys the combined envelope
\[
\frac{1}{256n}
\le
\mathfrak R^{+}_{n,d,q}
\le
\min\left\{
1,\,
\overline C\,r_{n,d,q},\,
\frac4n+(1-q)^2
\right\}.
\]
The clipped observed-outcome contrast \(\widehat\tau^{\mathrm{HT}}_n\) of \cref{obj:def:complete-arrival-ht-estimator} satisfies
\[
\sup_{P\in\mathcal M^{+}_{n,d,q}}
E_P\!\left[
\left\{
\widehat\tau^{\mathrm{HT}}_n(\mathbf O_n)-\tau(P)
\right\}^{2}
\right]
\le
\frac4n+(1-q)^2,
\]
where each expectation is taken under the canonical product law of \(n\) observed records induced by \(P\).
\end{theoremv}

\begin{theoremv}{thm:unrestricted-large-alphabet-deficit-lower}[Unrestricted risk with large alphabets]
Let \(n,d\) be nonnegative integers and \(q\) a real number satisfying:
\begin{itemize}
\item \textbf{(Sample size.)} \(n\ge 1\).
\item \textbf{(Alphabet size.)} \(d\ge 1024n^2\).
\item \textbf{(Arrival floor.)} \(0<q\le 1\).
\end{itemize}
Then \(\mathfrak R^{+}_{n,d,q}\), the unrestricted minimax squared-error risk defined in \cref{obj:def:unrestricted-minimax-risk}, satisfies
\[
\max\left\{\frac{1}{256n},\frac{(1-q)^2}{1024}\right\}
\le \mathfrak R^{+}_{n,d,q}
\le \frac{4}{n}+(1-q)^2.
\]
These bounds also give the envelope
\(\frac{1}{2048}\left\{\frac{1}{n}+(1-q)^2\right\}
\le \mathfrak R^{+}_{n,d,q}
\le 4\left\{\frac{1}{n}+(1-q)^2\right\}\).
\end{theoremv}

\begin{propositionv}{prop:unrestricted-uniform-near-complete-elbow}[Uniform near-complete arrival elbow]
Let \((q_n)_{n\ge0}\) satisfy \(0<q_n\le1\) for every \(n\), and let \(\mathfrak R^{+}_{n,d,q}\) be the unrestricted-arrival minimax squared-error risk defined in \cref{obj:def:unrestricted-minimax-risk}. For every integer \(n\ge1\),
\[
n\mathfrak R^{+}_{n,d,q_n}\le 4+n(1-q_n)^2
\qquad\text{for every integer }d\ge1.
\]
Moreover, for every real number
\[
b<\frac{1}{2048}\bigl\{1+n(1-q_n)^2\bigr\},
\]
there exists an integer \(d\ge1\) such that
\[
b<n\mathfrak R^{+}_{n,d,q_n}.
\]

The following two conditions are equivalent:
\begin{itemize}
\item \textbf{(Arrival rate.)} There exists a real constant \(M\ge0\) such that, for all sufficiently large integers \(n\),
\[
1-q_n\le \frac{M}{\sqrt n}.
\]
\item \textbf{(Uniform risk bound.)} There exists a real constant \(D\) such that, for all sufficiently large integers \(n\) and every integer \(d\ge1\),
\[
\mathfrak R^{+}_{n,d,q_n}\le \frac{D}{n}.
\]
\end{itemize}

For every pair of integers \(n,d\ge1\),
\[
\frac{1}{256n}\le \mathfrak R^{+}_{n,d,q_n}.
\]
If the sequence \(\bigl(\sqrt n(1-q_n)\bigr)_{n\ge0}\) is unbounded above, then for every real constant \(C\), there are infinitely many integers \(n\) such that
\[
C<n\mathfrak R^{+}_{n,\,1024n^2,\,q_n}.
\]
\end{propositionv}

\begin{propositionv}{prop:unrestricted-near-complete-phase-boundary}[Near-complete arrival phase boundary]
Let \(\mathfrak R^{+}_{n,d,q}\) be the unrestricted minimax squared-error risk defined in \cref{obj:def:unrestricted-minimax-risk}. Fix a real constant \(M\) and sequences \((d_n)_{n\ge0}\) and \((q_n)_{n\ge0}\) satisfying:
\begin{itemize}
\item \textbf{(Constant.)} \(M\ge0\).
\item \textbf{(Alphabet size.)} \(d_n\) is a natural number with \(d_n\ge1\) for every \(n\).
\item \textbf{(Arrival floor.)} \(q_n\in(0,1]\) for every \(n\).
\item \textbf{(Arrival deficit.)} For all sufficiently large \(n\),
\[
1-q_n\le \frac{M}{\sqrt n}.
\]
\end{itemize}
Then, for all sufficiently large \(n\),
\[
\frac{1}{256n}
\le \mathfrak R^{+}_{n,d_n,q_n}
\le \frac{4+M^2}{n}.
\]
\end{propositionv}

\begin{theoremv}{thm:unrestricted-complete-arrival-frontier}[Complete-arrival minimax risk frontier]
For every sample size \(n\ge 1\) and alphabet size \(d\ge 1\), let \(\mathcal M^{+}_{n,d,1}\) be the complete-arrival model class of \cref{obj:def:unrestricted-arrival-model-class}, and let \(\mathbf O_n\) consist of \(n\) independent, identically distributed observed records induced by the full-data law \(P\). Let \(\mathfrak R^{+}_{n,d,1}\) be the minimax squared-error risk of \cref{obj:def:unrestricted-minimax-risk}, and let \(\widehat\tau^{\mathrm{HT}}_n\) be the clipped observed-outcome arm contrast of \cref{obj:def:complete-arrival-ht-estimator}. Writing \(\tau(P)\) for the population average treatment effect, we have
\[
\frac{1}{256n}
\le \mathfrak R^{+}_{n,d,1}
\le
\sup_{P\in\mathcal M^{+}_{n,d,1}}
E_P\!\left[
\left\{\widehat\tau^{\mathrm{HT}}_n(\mathbf O_n)-\tau(P)\right\}^{2}
\right]
\le \frac{4}{n}.
\]
Consequently, at the complete-arrival endpoint \(q=1\), the minimax squared-error risk has order \(n^{-1}\) with universal constants, uniformly over \(d\ge 1\).
\end{theoremv}

\begin{propositionv}{prop:unrestricted-complete-arrival-phase-boundary}[Complete-arrival point-risk boundary]
There exist real constants \(0<c\le C\) such that, for every sequence of alphabet sizes \(d_n\ge1\), the unrestricted minimax squared-error risk \(\mathfrak R^{+}_{n,d_n,1}\) defined in \cref{obj:def:unrestricted-minimax-risk} at complete arrival \(q=1\) satisfies, for all sufficiently large \(n\),
\[
\frac{c}{n}\le \mathfrak R^{+}_{n,d_n,1}\le \frac{C}{n}.
\]
The comparison constants are independent of the chosen alphabet-size sequence.
\end{propositionv}

\begin{theoremv}{thm:matched-minimax-frontier}[Matched minimax frontier]
There exist universal constants \(0<\underline c\le\overline C<\infty\) such that, for every \(n,d\in\mathbb N\) and \(q\in\mathbb R\) satisfying
\begin{itemize}
\item \textbf{(Sample size and dimension.)} \(n\ge1\) and \(d\ge1\).
\item \textbf{(Arrival parameter.)} \(0<q\le1\).
\item \textbf{(Rare-arrival restriction.)} \cref{obj:ass:rare-arrival-slice} holds.
\end{itemize}
the minimax squared-error risk \(\mathfrak R_{n,d,q}\) defined in \cref{obj:def:minimax-risk} and the frontier rate \(r_{n,d,q}\) defined in \cref{obj:def:frontier-rate} satisfy
\[
\underline c\,r_{n,d,q}
\le \mathfrak R_{n,d,q}
\le \overline C\,r_{n,d,q}.
\]
\end{theoremv}

\begin{propositionv}{prop:fixed-d-effective-sample-reduction}[Fixed-alphabet effective sample reduction]
Fix an integer \(d\ge 1\). Let \(n(t)\) and \(q(t)\) be sequences indexed by \(t\in\mathbb N\) satisfying:
\begin{itemize}
\item \textbf{(Sample size and arrival.)} For every \(t\), \(n(t)\in\mathbb N\), \(n(t)\ge 1\), and \(0<q(t)\le 1\).
\item \textbf{(Rare-arrival slice.)} For every \(t\), the pair \((n(t),q(t))\) satisfies \cref{obj:ass:rare-arrival-slice}.
\item \textbf{(Effective sample size.)} The effective sample size
\[
N(t)=n(t)q(t)
\]
satisfies \(N(t)\to\infty\) as \(t\to\infty\).
\end{itemize}
Then the frontier rate and minimax squared-error risk of \cref{obj:def:frontier-rate,obj:def:minimax-risk} satisfy
\[
r_{n(t),d,q(t)}\asymp N(t)^{-1},
\qquad
\mathfrak R_{n(t),d,q(t)}\asymp N(t)^{-1}
\qquad (t\to\infty).
\]
The minimax squared-error risk with the same estimator class and loss as in \cref{obj:def:minimax-risk}, taken over the restricted full-data laws
\[
\bigl\{P\in\mathcal M_{n(t),d,q(t)}:
P\bigl(S(0)=0,\ S(1)=0\bigr)=1\bigr\},
\]
is also of order \(N(t)^{-1}\) as \(t\to\infty\). Each order comparison means eventual upper and lower bounds by positive constant multiples of \(N(t)^{-1}\).
\end{propositionv}

\begin{propositionv}{prop:phase-boundaries}[Rare-arrival minimax phase boundaries]
Let \((n_t)_{t\in\mathbb N}\) and \((d_t)_{t\in\mathbb N}\) be sequences of natural numbers, and let \((q_t)_{t\in\mathbb N}\) be a sequence of real numbers. Define the effective sample size by
\[
N_t=n_tq_t.
\]
Suppose that:
\begin{itemize}
\item \textbf{(Positive sizes.)} \(n_t\ge 1\) and \(d_t\ge 1\) for every \(t\in\mathbb N\).
\item \textbf{(Arrival probabilities.)} \(0<q_t\le 1\) for every \(t\in\mathbb N\).
\item \textbf{(Rare-arrival slice.)} \((n_t,q_t)\) satisfies \cref{obj:ass:rare-arrival-slice} for every \(t\in\mathbb N\).
\item \textbf{(Effective sample growth.)} \(N_t\to\infty\) as \(t\to\infty\).
\end{itemize}
Then the minimax squared-error risk \(\mathfrak R_{n_t,d_t,q_t}\) of \cref{obj:def:minimax-risk} satisfies both equivalences:
\[
\mathfrak R_{n_t,d_t,q_t}\longrightarrow 0
\quad\Longleftrightarrow\quad
d_t=o\!\left(N_t\log(e+N_t)\right),
\]
\[
\mathfrak R_{n_t,d_t,q_t}=\Theta\!\left(N_t^{-1}\right)
\quad\Longleftrightarrow\quad
d_t=O\!\left(\sqrt{N_t}\log(e+N_t)\right),
\]
where all limits and asymptotic comparisons are taken as \(t\to\infty\).
\end{propositionv}

\begin{theoremv}{thm:unrestricted-fixed-q-minimax-frontier}[Unrestricted fixed-arrival-floor minimax frontier]*
Fix an arrival floor \(q\) with \(0<q<1/2\). There exist real constants \(0<\underline c(q)\le \overline C(q)<\infty\), depending only on \(q\), such that for every pair of integers \(n,d\ge1\), the unrestricted minimax squared-error risk \(\mathfrak R^{+}_{n,d,q}\) defined in \cref{obj:def:unrestricted-minimax-risk} satisfies
\[
\underline c(q)
\min\left\{1,\frac1n+\left(\frac{d}{n\log(e+n)}\right)^2\right\}
\le \mathfrak R^{+}_{n,d,q}
\le \overline C(q)
\min\left\{1,\frac1n+\left(\frac{d}{n\log(e+n)}\right)^2\right\}.
\]
With the same constants, the frontier rate \(r_{n,d,q}\) defined in \cref{obj:def:frontier-rate} satisfies
\[
\underline c(q)\,r_{n,d,q}
\le \mathfrak R^{+}_{n,d,q}
\le \overline C(q)\,r_{n,d,q}.
\]
\end{theoremv}
% CAUSALSMITH-CITED-SCOPE-BEGIN thm:unrestricted-fixed-q-minimax-frontier
\verificationfootnotetext{\textbf{Formalization scope.} This result uses the cited conclusion \citep{ZengBalakrishnanHanKennedy2026} (Theorem 2, Section 3.2, p. 9, together with its proof in Section C.5, pp. 39--44, which sets the control regression to zero). The cited source proof is not formalized here: Lean verifies its use and all remaining steps. If that cited conclusion is false, the portions of this result that depend on it are not certified.}
% CAUSALSMITH-CITED-SCOPE-END thm:unrestricted-fixed-q-minimax-frontier

\begin{propositionv}{prop:unrestricted-fixed-q-phase-boundaries}[Unrestricted fixed-arrival-floor phase boundaries]*
Fix \(0<q<1/2\), where \(q\) is the arrival floor, and let \(\mathfrak R^{+}_{n,d,q}\) be the unrestricted-arrival minimax squared-error risk defined in \cref{obj:def:unrestricted-minimax-risk}. For every sequence of integers \(d_n\ge1\),
\[
\mathfrak R^{+}_{n,d_n,q}\longrightarrow0
\quad\Longleftrightarrow\quad
d_n=o\!\left(n\log(e+n)\right),
\qquad n\to\infty.
\]

There exists a real constant \(c>0\) such that, for every real \(M>0\), there exists a real constant \(C\ge c\) with the following property: for every sequence of integers \(d_n\ge1\), if
\[
d_n\le M\sqrt n\,\log(e+n)
\qquad\text{for all sufficiently large }n,
\]
then
\[
\frac{c}{n}\le\mathfrak R^{+}_{n,d_n,q}\le\frac{C}{n}
\qquad\text{for all sufficiently large }n.
\]
Here \(c\) may depend on the fixed \(q\), and \(C\) may depend on \(q\) and \(M\); the thresholds for sufficiently large \(n\) may depend on the sequence.

Conversely, for every sequence of integers \(d_n\ge1\), if there exist real constants \(0<c_*\le C_*\) such that
\[
\frac{c_*}{n}\le\mathfrak R^{+}_{n,d_n,q}\le\frac{C_*}{n}
\qquad\text{for all sufficiently large }n,
\]
then
\[
d_n=O\!\left(\sqrt n\,\log(e+n)\right),
\qquad n\to\infty.
\]
\end{propositionv}
% CAUSALSMITH-CITED-SCOPE-BEGIN prop:unrestricted-fixed-q-phase-boundaries
\verificationfootnotetext{\textbf{Formalization scope.} This result uses the cited conclusion \citep{ZengBalakrishnanHanKennedy2026} (Theorem 2, Section 3.2, p. 9, together with its proof in Section C.5, pp. 39--44, which sets the control regression to zero). The cited source proof is not formalized here: Lean verifies its use and all remaining steps. If that cited conclusion is false, the portions of this result that depend on it are not certified.}
% CAUSALSMITH-CITED-SCOPE-END prop:unrestricted-fixed-q-phase-boundaries

\begin{theoremv}{thm:unrestricted-mixed-count-upper}[Unrestricted mixed-count risk bound]
There exists a universal constant \(0<\overline C<\infty\) such that, for every sample size \(n\), alphabet size \(d\), arrival floor \(q\), and full-data law \(P\) satisfying
\begin{itemize}
\item \textbf{(Domain.)} \(n,d\ge1\) and \(0<q\le1\).
\item \textbf{(Model.)} \(P\in\mathcal M^{+}_{n,d,q}\), the model class defined in \cref{obj:def:unrestricted-arrival-model-class}.
\item \textbf{(Sampling.)} The observed sample \(\mathbf O_n\) has the canonical product law induced by \(P\), with its coordinate observations satisfying \cref{obj:ass:iid-sampling}.
\end{itemize}
the deterministic estimator \(\widehat\tau^{\mathrm{mix}}_{n,d,q}\) of \cref{obj:def:mixed-count-estimator} satisfies
\[
E_P\!\left[
 \bigl(\widehat\tau^{\mathrm{mix}}_{n,d,q}(\mathbf O_n)-\tau(P)\bigr)^2
\right]
\le \mathfrak U_{n,d,q}
\le \overline C\,r_{n,d,q},
\]
where \(\tau(P)\) is the average treatment effect defined in \cref{obj:def:ate-functional}, \(r_{n,d,q}\) is the rate defined in \cref{obj:def:frontier-rate}, and \(\mathfrak U_{n,d,q}\) is the deterministic risk envelope specified below.

Set
\[
N=nq,\qquad m=\frac n6,\qquad N_0=mq=\frac N6,
\qquad \gamma_0=\log 2-\frac12.
\]
Use the logarithmic scale \(\ell\) defined in \cref{obj:synth_7}, the polynomial degree \(k\) defined in \cref{obj:synth_8}, and the threshold \(B\) defined in \cref{obj:synth_18}. Define the certificate quantities
\[
b_{n,d,q}
=2\left(\frac{4dB}{N_0k^2}+3e^{-16\ell}\right),
\]
\[
v_{n,d,q}
=8\left(\frac4{N_0}+\frac1m\right)
+64d\left(e^{\ell/8}+1\right)
 \left(\frac{B^2}{N_0^2}+\frac{B}{mN_0}\right)
+32e^{-32\ell}\left(1+\frac1m\right).
\]
For \(N\le1\), set \(\mathfrak U_{n,d,q}=1\). For \(N>1\) on the elementary branch of \cref{obj:def:mixed-count-estimator}, set
\[
\mathfrak U_{n,d,q}
=\min\left\{
4,\,
\left(\frac{8d}{eN_0}\right)^2
+4\left(\frac4{N_0}+\frac1m\right)
+4e^{-n\gamma_0}
\right\}.
\]
For \(N>1\) on its polynomial branch, set
\[
\mathfrak U_{n,d,q}
=\min\left\{
4,\,
b_{n,d,q}^{\,2}+v_{n,d,q}+4e^{-n\gamma_0}
\right\}.
\]

Moreover, the estimator's conditional average over the auxiliary randomization satisfies
\[
E_P\!\left[
 \bigl(\widehat\tau^{\mathrm{mix}}_{n,d,q}(\mathbf O_n)-\tau(P)\bigr)^2
\right]
\le
E_P\!\left[
 E_{\omega}\!\left[
  \bigl(\widetilde\tau(\mathbf O_n;\omega)-\tau(P)\bigr)^2
  \,\middle|\,\mathbf O_n
 \right]
\right],
\]
where \(\omega\) is the auxiliary prefix and stream-assignment randomization from \cref{obj:def:mixed-count-estimator}, and \(\widetilde\tau(\mathbf O_n;\omega)\) is its output, including the zero output when the Poisson prefix length exceeds \(n\).
\end{theoremv}

\begin{definitionv}{def:interval-length-risk}[Minimax expected interval length]
The minimax maximal expected interval length, \(\mathfrak L_{n,d,q,\alpha}\), is defined for \(n,d\in\mathbb N\) and \(q,\alpha\in\mathbb R\) as follows. Let \(\mathcal M_{n,d,q}\) be the model class in \cref{obj:def:model-class}, let \(\mathbf O_n\) be the observed sample with law \(\mathcal L_P(\mathbf O_n)\) from \cref{obj:def:observed-experiment}, and let \(\tau(P)\) be the average treatment effect in \cref{obj:def:ate-functional}.

The class \(\mathcal C_{n,d,q,\alpha}\) consists of parameter-independent Markov kernels \(\mathsf I\) from the observed-sample space to endpoint pairs \((l,u)\) satisfying \(-1\le l\le u\le1\), with uniform joint sample-and-randomization coverage: for every \(P\in\mathcal M_{n,d,q}\),
\[
\int\!\int
\mathbf1\{l\le\tau(P)\le u\}\,
\mathsf I(o,d(l,u))\,
\mathcal L_P(\mathbf O_n)(do)
\ge 1-\alpha.
\]
The realized connected interval is \([l,u]\), with length \(u-l\). Define
\[
\mathfrak L_{n,d,q,\alpha}
:=
\inf_{\mathsf I\in\mathcal C_{n,d,q,\alpha}}
\sup_{P\in\mathcal M_{n,d,q}}
\int\!\int (u-l)\,
\mathsf I(o,d(l,u))\,
\mathcal L_P(\mathbf O_n)(do),
\]
where the infimum and supremum are real-valued, with the convention that an infimum or supremum over an empty set equals zero. In particular, \(\mathfrak L_{n,d,q,\alpha}=0\) when \(\mathcal C_{n,d,q,\alpha}\) is empty. Both coverage and expected length integrate over the sample and the kernel draw. Deterministic connected-interval procedures are included through Dirac kernels.
\end{definitionv}

\begin{definitionv}{def:risk-envelope-interval}[Risk-envelope interval \(I^{\mathrm{mix}}_\alpha\)]
The risk-envelope interval is
\[
I^{\mathrm{mix}}_\alpha
:=
\left[
\widehat\tau^{\mathrm{mix}}_{n,d,q}
-
\sqrt{\frac{\mathfrak U_{n,d,q}}{\alpha}},
\,
\widehat\tau^{\mathrm{mix}}_{n,d,q}
+
\sqrt{\frac{\mathfrak U_{n,d,q}}{\alpha}}
\right]\cap[-1,1],
\]
where \(\widehat\tau^{\mathrm{mix}}_{n,d,q}\) is the mixed-count estimator and \(\mathfrak U_{n,d,q}\) is its explicit deterministic squared-error envelope.
\end{definitionv}

\begin{theoremv}{thm:connected-interval-frontier}[Connected interval length frontier]
For every miscoverage level \(\alpha\in(0,1)\), there exist constants \(0<\underline c_\alpha\le\overline C_\alpha<\infty\) such that the following holds for every \(n,d\in\mathbb N\) and \(q\in\mathbb R\) satisfying:
\begin{itemize}
\item \textbf{(Sample size and dimension.)} \(n\ge1\) and \(d\ge1\).
\item \textbf{(Arrival floor.)} \(0<q\le1\).
\item \textbf{(Rare-arrival slice.)} The restriction in \cref{obj:ass:rare-arrival-slice} holds.
\end{itemize}
Let \(r_{n,d,q}\) be the frontier rate in \cref{obj:def:frontier-rate}, let \(\mathfrak L_{n,d,q,\alpha}\) be the minimax expected interval length in \cref{obj:def:interval-length-risk}, and let \(I^{\mathrm{mix}}_\alpha\) be the interval in \cref{obj:def:risk-envelope-interval}. Over the model class \(\mathcal M_{n,d,q}\) in \cref{obj:def:model-class},
\[
\underline c_\alpha\sqrt{r_{n,d,q}}
\le \mathfrak L_{n,d,q,\alpha}
\le \sup_{P\in\mathcal M_{n,d,q}}
E_P\!\left[\operatorname{length}(I^{\mathrm{mix}}_\alpha)\right]
\le \overline C_\alpha\sqrt{r_{n,d,q}}.
\]
Here expectations and probabilities are taken under the law of the independent observed sample induced by \(P\), and interval length is defined by
\[
\operatorname{length}([l,u])=u-l.
\]
Moreover, for every \(P\in\mathcal M_{n,d,q}\), the average treatment effect \(\tau(P)\) from \cref{obj:def:ate-functional} satisfies
\[
\Pr_P\!\left\{\tau(P)\in I^{\mathrm{mix}}_\alpha\right\}\ge1-\alpha.
\]
\end{theoremv}

\begin{definitionv}{def:zeng-discrete-ate-class}[Observational law class \(\mathcal D_{d,\epsilon}\)]
The class \(\mathcal D_{d,\epsilon}\), for \(0<\epsilon<1/2\), consists of one-record laws \(P_Z\) of the baseline category, observational treatment, and outcome
\[
(X,B^{\mathrm{obs}},Y)\in\{1,\ldots,d\}\times\{0,1\}^2
\]
with baseline probabilities
\[
p_x=P_Z(X=x)
\]
and, on each occupied category \(p_x>0\), conditional treatment probabilities and outcome means
\[
\pi_x=P_Z(B^{\mathrm{obs}}=1\mid X=x),
\qquad
\mu_{bx}=E_{P_Z}[Y\mid X=x,B^{\mathrm{obs}}=b],
\]
satisfying
\[
\epsilon\le\pi_x\le1-\epsilon,
\qquad
0\le\mu_{bx}\le1
\quad (b\in\{0,1\}).
\]
The conditional quantities \(\pi_x\) and \(\mu_{bx}\) are specified on occupied categories.
\end{definitionv}

\begin{definitionv}{def:zeng-ate-functional}[Observed-law contrast $\theta_Z$]
The total observed-law functional \(\theta_Z\) is defined, for every \(P_Z\in\mathcal D_{d,\epsilon}\), by
\[
\theta_Z(P_Z)
:=
\sum_{x:p_x>0}p_x(\mu_{1x}-\mu_{0x}).
\]
Null baseline categories contribute zero.
\end{definitionv}

\begin{definitionv}{def:zeng-minimax-risk}[Observational minimax risk \(\mathfrak R^Z_{n,d,\epsilon}\)]
The minimax squared-error risk for the observational contrast is
\[
\mathfrak R^Z_{n,d,\epsilon}
:=\inf_{\widehat\theta}
\sup_{P_Z\in\mathcal D_{d,\epsilon}}
E_{P_Z}\!\left[
(\widehat\theta(\mathbf Z_n)-\theta_Z(P_Z))^2
\right],
\]
where \(\mathbf Z_n\) denotes the independent, identically distributed sample of \(n\) observational records with law \(P_Z\). The infimum ranges over all possibly randomized estimators measurable with respect to \(\mathbf Z_n\) and auxiliary randomness independent of \(P_Z\); the expectation integrates both sources of randomness.
\end{definitionv}

\begin{theoremv}{thm:zeng-fixed-positivity-polynomial-frontier}[Fixed positivity polynomial frontier]*
There exists a universal constant \(\overline C>0\) such that the following holds for every pair of integers \(n,d\ge1\) and every strict positivity margin \(0<\epsilon<1/2\). Let \(\mathcal D_{d,\epsilon}\) be the observational class in \cref{obj:def:zeng-discrete-ate-class}, and let \(\mathbf Z_n=((X_i,B_i^{\mathrm{obs}},Y_i))_{i=1}^n\) be an independent sample from \(P_Z\in\mathcal D_{d,\epsilon}\). Write the observational baseline categories as \(1,\ldots,d\). For a realization \(\mathbf Z_n=((x_i,b_i,y_i))_{i=1}^n\), define the synthetic records, with baseline coordinates indexed from zero, by
\[
o_i(u_i)
=
\bigl(x_i-1,u_i,0,1\{b_i=u_i\},1\{b_i=u_i\}y_i\bigr),
\qquad u_i\in\{0,1\},
\]
and define the deterministic estimator
\[
\widehat\theta^{\mathrm{poly}}_{n,d,\epsilon}(\mathbf Z_n)
=
\frac{1}{2^n}
\sum_{(u_1,\ldots,u_n)\in\{0,1\}^n}
\widehat\tau^{\mathrm{mix}}_{n,d,\epsilon}
\bigl(o_1(u_1),\ldots,o_n(u_n)\bigr),
\]
where \(\widehat\tau^{\mathrm{mix}}_{n,d,\epsilon}\) is the mixed-count estimator in \cref{obj:def:mixed-count-estimator}. Then
\[
\sup_{P_Z\in\mathcal D_{d,\epsilon}}
E_{P_Z}\!\left[
\left(
\widehat\theta^{\mathrm{poly}}_{n,d,\epsilon}(\mathbf Z_n)
-\theta_Z(P_Z)
\right)^2
\right]
\le
\overline C\,r_{n,d,\epsilon},
\]
where \(\theta_Z(P_Z)\) is the observational contrast in \cref{obj:def:zeng-ate-functional} and \(r_{n,d,\epsilon}\) is the rate in \cref{obj:def:frontier-rate}.

Moreover, for each fixed \(0<\epsilon<1/2\), there exist constants \(0<c_\epsilon\le C_\epsilon<\infty\) such that, for every pair of integers \(n,d\ge1\), the minimax risk in \cref{obj:def:zeng-minimax-risk} satisfies
\[
c_\epsilon
\min\left\{
1,\frac1n+
\left[\frac{d}{n\log(e+n)}\right]^2
\right\}
\le
\mathfrak R^Z_{n,d,\epsilon}
\le
C_\epsilon
\min\left\{
1,\frac1n+
\left[\frac{d}{n\log(e+n)}\right]^2
\right\}.
\]
\end{theoremv}
% CAUSALSMITH-CITED-SCOPE-BEGIN thm:zeng-fixed-positivity-polynomial-frontier
\verificationfootnotetext{\textbf{Formalization scope.} This result uses the cited conclusion \citep{ZengBalakrishnanHanKennedy2026} (Theorem 2, Section 3.2, p. 9, together with its proof in Section C.5, pp. 39--44, which sets the control regression to zero). The cited source proof is not formalized here: Lean verifies its use and all remaining steps. If that cited conclusion is false, the portions of this result that depend on it are not certified.}
% CAUSALSMITH-CITED-SCOPE-END thm:zeng-fixed-positivity-polynomial-frontier

\begin{lemmav}{lem:randomized-kernel-barycenter}[Barycenters and uniform randomization]
Let \(\Lambda\) be an arbitrary parameter set, and suppose:
\begin{itemize}
\item \textbf{(Sample space.)} \(\mathcal Z\) is a finite measurable space with measurable singletons.
\item \textbf{(Experiment.)} For each \(\lambda\in\Lambda\), \(Q_\lambda\) is a probability law on \(\mathcal Z\).
\item \textbf{(Target.)} For each \(\lambda\in\Lambda\), the target satisfies \(\vartheta(\lambda)\in[-1,1]\).
\end{itemize}
For a parameter-independent Markov kernel \(\kappa\) from \(\mathcal Z\) to \(\mathbb R\), supported on \([-1,1]\) at every observation, define its barycenter by
\[
b_\kappa(o)=\int_{\mathbb R}h\,\kappa(o,dh).
\]
Then \(b_\kappa(o)\in[-1,1]\) for every \(o\in\mathcal Z\), and, for every \(\lambda\in\Lambda\),
\[
\int_{\mathcal Z}\bigl(b_\kappa(o)-\vartheta(\lambda)\bigr)^2\,Q_\lambda(do)
\le
\int_{\mathcal Z}\int_{\mathbb R}
\bigl(h-\vartheta(\lambda)\bigr)^2\,\kappa(o,dh)\,Q_\lambda(do).
\]
The minimax squared risk over these bounded kernels equals that over deterministic maps into \([-1,1]\):
\[
\inf_{\kappa}\sup_{\lambda\in\Lambda}
\int_{\mathcal Z}\int_{\mathbb R}
\bigl(h-\vartheta(\lambda)\bigr)^2\,\kappa(o,dh)\,Q_\lambda(do)
=
\inf_{b:\mathcal Z\to[-1,1]}\sup_{\lambda\in\Lambda}
\int_{\mathcal Z}\bigl(b(o)-\vartheta(\lambda)\bigr)^2\,Q_\lambda(do),
\]
where the first infimum ranges over the bounded kernels just specified.

Define clipping by
\[
\operatorname{clip}_{[-1,1]}(h)=\max\{-1,\min\{1,h\}\}.
\]
For every parameter-independent Markov kernel \(\kappa\) from \(\mathcal Z\) to \(\mathbb R\), and every \(\lambda\in\Lambda\), its image under clipping satisfies
\[
\int_{\mathcal Z}\int_{\mathbb R}
\bigl(\operatorname{clip}_{[-1,1]}(h)-\vartheta(\lambda)\bigr)^2
\,\kappa(o,dh)\,Q_\lambda(do)
\le
\int_{\mathcal Z}\int_{\mathbb R}
\bigl(h-\vartheta(\lambda)\bigr)^2\,\kappa(o,dh)\,Q_\lambda(do),
\]
with both risks interpreted as nonnegative extended integrals. The extended minimax squared risk over all such real-valued Markov kernels, converted to a real value, equals the bounded-kernel minimax squared risk above.

The uniform law on \([0,1]\), defined by restricting Lebesgue measure to this interval, is a probability law. Every bounded kernel \(\kappa\) admits a measurable map
\(f:\mathcal Z\times[0,1]\to[-1,1]\) such that, for every \(o\in\mathcal Z\), \(\kappa(o,\cdot)\) is the distribution of \(f(o,U)\), where \(U\) has this uniform law. Conversely, every such measurable map induces a bounded Markov kernel with precisely these conditional distributions. The same uniform seed law serves every parameter.
\end{lemmav}

\begin{lemmav}{lem:randomized-kernel-tv-comparison}[Total variation under randomization]
Let \(Q_0,Q_1\) be probability laws on a finite measurable space \(\mathcal Z\) with measurable singletons, and let \(\kappa\) be a common Markov kernel from \(\mathcal Z\) to a measurable action space \(\mathcal H\). For \(i\in\{0,1\}\), write \(Q_i\ltimes\kappa\) for the joint law obtained by drawing an input from \(Q_i\) and then an action from \(\kappa\) conditional on that input, and write \(Q_i\kappa\) for its action marginal. For two probability laws on the same measurable space, define their total variation distance as the supremum of their absolute probability differences over measurable events. Then
\[
\operatorname{TV}(Q_0\kappa,Q_1\kappa)
\le
\operatorname{TV}(Q_0\ltimes\kappa,Q_1\ltimes\kappa)
=
\operatorname{TV}(Q_0,Q_1).
\]
For every measurable event \(T\subseteq\mathcal Z\times\mathcal H\),
\[
(Q_0\ltimes\kappa)(T)+(Q_1\ltimes\kappa)(T^c)
\ge 1-\operatorname{TV}(Q_0,Q_1),
\]
and for every measurable event \(T\subseteq\mathcal H\),
\[
(Q_0\kappa)(T)+(Q_1\kappa)(T^c)
\ge 1-\operatorname{TV}(Q_0,Q_1).
\]
\end{lemmav}

\begin{lemmav}{lem:unrestricted-cell-identification}[Identification by cell contributions]
Let \(n,d\in\mathbb N\), \(q\in\mathbb R\), and let \(P\in\mathcal M^{+}_{n,d,q}\) satisfy \cref{obj:def:unrestricted-arrival-model-class}. The average treatment effect \(\tau(P)\) and the total observed-cell functional \(\Phi(P)\), defined in \cref{obj:def:ate-functional,obj:def:cell-functional}, satisfy
\[
\tau(P)=\Phi(P)
=2\sum_{(a,x,s)\in\mathcal J_d}g(a)c_{axs}(P).
\]
\end{lemmav}

\begin{lemmav}{lem:cell-identification-background}[Identification by cell contributions]
Let \(n,d\in\mathbb N\), \(q\in\mathbb R\), and let \(P\in\mathcal M_{n,d,q}\) satisfy the model conditions in \cref{obj:def:model-class}. The average treatment effect \(\tau(P)\) of \cref{obj:def:ate-functional} equals the total signed observed-cell functional \(\Phi(P)\) of \cref{obj:def:cell-functional}:
\[
\tau(P)=\Phi(P)
=2\sum_{(a,x,s)\in\mathcal J_d}g(a)c_{axs}(P).
\]
Here the cell contributions and their null-cell convention are as defined in \cref{obj:def:cell-functional}.
\end{lemmav}

\begin{lemmav}{lem:synthetic-randomization-kernel}[Synthetic randomization kernel]
Let the inputs satisfy:
\begin{itemize}
\item \textbf{(Sample size.)} \(n,d\in\mathbb N\) and \(n\ge1\).
\item \textbf{(Arrival floor.)} \(0<q<1/2\).
\item \textbf{(Observational law.)} \(P_Z\in\mathcal D_{d,q}\), the discrete positivity class of \cref{obj:def:zeng-discrete-ate-class}.
\end{itemize}
Index observational baseline categories by \(x\in\{1,\ldots,d\}\) and full-data baseline categories by their corresponding zero-based indices \(x-1\in\{0,\ldots,d-1\}\). For an observational record \((x,b,y)\), define
\[
o(x,b,y;u)
=
\left(x-1,u,0,\mathbf 1\{b=u\},\mathbf 1\{b=u\}y\right),
\qquad u\in\{0,1\}.
\]
Applying this map independently to each input record with independent fair binary draws defines a Markov kernel independent of \(P_Z\).

There exists a full-data law \(P_Z^*\in\mathcal M^{+}_{n,d,q}\), as defined in \cref{obj:def:unrestricted-arrival-model-class}, such that applying this kernel to \(\mathbf Z_n\), an i.i.d.\ sample of \(n\) observational records from \(P_Z\), yields exactly the \(n\)-fold product of the observed-record law induced by \(P_Z^*\). The resulting observed sample \(\mathbf O_n\) satisfies \cref{obj:ass:iid-sampling} under \(P_Z^*\), and
\[
\tau(P_Z^*)=\theta_Z(P_Z),
\]
with the functionals defined in \cref{obj:def:ate-functional,obj:def:zeng-ate-functional}.
\end{lemmav}

\begin{theoremv}{thm:unrestricted-uniform-mixed-count-certificate}[Uniform mixed-count risk certificate]
There exists a universal constant \(0<\overline C<\infty\) such that the following holds.
\begin{itemize}
\item \textbf{(Domain.)} Let \(n,d\ge1\) be integers and let \(0<q\le1\).
\item \textbf{(Model.)} Let \(P\in\mathcal M^{+}_{n,d,q}\), the unrestricted-arrival model class of \cref{obj:def:unrestricted-arrival-model-class}.
\item \textbf{(Sampling.)} Let \(\mathbf O_n\), the sample of observed records, satisfy \cref{obj:ass:iid-sampling} under \(P\).
\end{itemize}
Use the mixed-count estimator of \cref{obj:def:mixed-count-estimator}. Let \(\ell\), the logarithmic scale, be defined by \cref{obj:synth_7}; let \(k\), the polynomial degree, be defined by \cref{obj:synth_8}; and let \(B\), the threshold, be defined by \cref{obj:synth_18}. Put
\[
N=nq,\qquad m=\frac n6,\qquad N_0=mq=\frac N6,
\qquad \gamma_0=\log 2-\frac12.
\]
Define the deterministic envelope \(\mathfrak U_{n,d,q}\) as follows. If \(N\le1\), set
\[
\mathfrak U_{n,d,q}=1.
\]
If \(N>1\) and either \(\ell<128\) or \(d<\sqrt N\), set
\[
\mathfrak U_{n,d,q}
=\min\left\{4,
\left(\frac{8d}{eN_0}\right)^2
+4\left(\frac4{N_0}+\frac1m\right)
+4e^{-n\gamma_0}\right\}.
\]
If \(N>1\), \(\ell\ge128\), and \(d\ge\sqrt N\), put
\[
b_{n,d,q}
=2\left\{\frac{4dB}{N_0k^2}+3e^{-16\ell}\right\},
\]
\[
v_{n,d,q}
=8\left(\frac4{N_0}+\frac1m\right)
+64d(e^{\ell/8}+1)
\left(\frac{B^2}{N_0^2}+\frac{B}{mN_0}\right)
+32e^{-32\ell}\left(1+\frac1m\right),
\]
and set
\[
\mathfrak U_{n,d,q}
=\min\left\{4,b_{n,d,q}^2+v_{n,d,q}
+4e^{-n\gamma_0}\right\}.
\]
Then, with \(\Phi(P)\) the cell functional of \cref{obj:def:cell-functional} and \(r_{n,d,q}\) the rate of \cref{obj:def:frontier-rate}, expectation under the canonical i.i.d. observed-sample law satisfies
\[
E_P\!\left[
\left(\widehat\tau^{\mathrm{mix}}_{n,d,q}(\mathbf O_n)-\Phi(P)\right)^2
\right]
\le \mathfrak U_{n,d,q}
\le \overline C\,r_{n,d,q}.
\]
\end{theoremv}

\begin{theoremv}{thm:uniform-mixed-count-certificate}[Uniform mixed-count risk certificate]
There exists a universal constant \(0<\overline C<\infty\) such that the following holds for every \(n,d\in\mathbb N\), \(q\in\mathbb R\), and full-data law \(P\in\mathcal M_{n,d,q}\) satisfying \cref{obj:def:model-class}, including \(n\ge1\), \(d\ge1\), and \(0<q\le1\).

Let \(\widehat\tau^{\mathrm{mix}}_{n,d,q}\) be the estimator in \cref{obj:def:mixed-count-estimator}, and let \(r_{n,d,q}\) be the rate in \cref{obj:def:frontier-rate}. Let \(\ell\) be the logarithmic scale defined in \cref{obj:synth_7}, \(k\) the polynomial degree defined in \cref{obj:synth_8}, and \(B\) the threshold defined in \cref{obj:synth_18}. Put
\[
N=nq,\qquad m=\frac n6,\qquad N_0=mq=\frac N6,\qquad
\gamma_0=\log 2-\frac12,
\]
where \(N\) is the effective sample size, \(m\) the stream intensity, \(N_0\) the effective stream size, and \(\gamma_0\) the prefix-tail constant. Define the deterministic envelope \(\mathfrak U_{n,d,q}\) by the following branches:
\begin{itemize}
\item \textbf{(Small effective sample.)} If \(N\le1\), set
\[
\mathfrak U_{n,d,q}=1.
\]
\item \textbf{(Elementary branch.)} If \(N>1\) and either \(\ell<128\) or \(d<\sqrt N\), set
\[
\mathfrak U_{n,d,q}
=\min\left\{4,
\left(\frac{8d}{eN_0}\right)^2
+4\left(\frac4{N_0}+\frac1m\right)
+4e^{-n\gamma_0}\right\}.
\]
\item \textbf{(Polynomial branch.)} If \(N>1\), \(\ell\ge128\), and \(d\ge\sqrt N\), define the certificate quantities
\[
b_{n,d,q}
=2\left(\frac{4dB}{N_0k^2}+3e^{-16\ell}\right),
\]
\[
v_{n,d,q}
=8\left(\frac4{N_0}+\frac1m\right)
+64d\left(e^{\ell/8}+1\right)
\left(\frac{B^2}{N_0^2}+\frac{B}{mN_0}\right)
+32e^{-32\ell}\left(1+\frac1m\right),
\]
and set
\[
\mathfrak U_{n,d,q}
=\min\left\{4,b_{n,d,q}^2+v_{n,d,q}+4e^{-n\gamma_0}\right\}.
\]
\end{itemize}
Then, writing \(\mathbf O_n\) for the observed sample induced by \(P\) and \(\tau(P)\) for the average treatment effect in \cref{obj:def:ate-functional},
\[
E_P\!\left[
\left(\widehat\tau^{\mathrm{mix}}_{n,d,q}(\mathbf O_n)-\tau(P)\right)^2
\right]
\le \mathfrak U_{n,d,q}
\le \overline C\,r_{n,d,q}.
\]
\end{theoremv}

\begin{theoremv}{thm:mixed-count-upper}[Mixed-count risk upper bound]
There exists a universal constant \(0<\overline C<\infty\) such that the following holds for every choice of parameters and law satisfying:
\begin{itemize}
\item \textbf{(Sample size and dimension.)} \(n,d\in\mathbb N\) satisfy \(n,d\ge1\).
\item \textbf{(Arrival floor.)} \(q\in\mathbb R\) satisfies \(0<q\le1\).
\item \textbf{(Model.)} The full-data law \(P\) belongs to \(\mathcal M_{n,d,q}\) in \cref{obj:def:model-class}.
\end{itemize}
Let \(\widehat\tau^{\mathrm{mix}}_{n,d,q}\) and \(\widetilde\tau(\mathbf O_n;\omega)\) be the deterministic estimator and auxiliary randomized output of \cref{obj:def:mixed-count-estimator}, where \(\mathbf O_n\) is the observed sample and \(\omega\) is the auxiliary randomization. Let \(\ell\) be the logarithmic scale defined in \cref{obj:synth_7}, \(k\) the polynomial degree defined in \cref{obj:synth_8}, and \(B=256\ell\) the threshold defined in \cref{obj:synth_18}. Use the elementary and polynomial branch rules of \cref{obj:def:mixed-count-estimator}. Define the effective sample size \(N\), stream intensity \(m\), effective stream size \(N_0\), and prefix-tail constant \(\gamma_0\) by
\[
N=nq,\qquad m=\frac n6,\qquad N_0=mq=\frac N6,
\qquad \gamma_0=\log 2-\frac12.
\]
Define the certificate quantities
\[
b_{n,d,q}
=
2\left(\frac{4dB}{N_0k^2}+3e^{-16\ell}\right)
\]
and
\[
v_{n,d,q}
=
8\left(\frac4{N_0}+\frac1m\right)
+
64d\left(e^{\ell/8}+1\right)
\left(\frac{B^2}{N_0^2}+\frac{B}{mN_0}\right)
+
32e^{-32\ell}\left(1+\frac1m\right).
\]
The deterministic risk envelope \(\mathfrak U_{n,d,q}\) is defined as follows. For \(N\le1\), set
\[
\mathfrak U_{n,d,q}=1.
\]
For \(N>1\) on the elementary branch, set
\[
\mathfrak U_{n,d,q}
=
\min\left\{
4,\,
\left(\frac{8d}{eN_0}\right)^2
+4\left(\frac4{N_0}+\frac1m\right)
+4e^{-n\gamma_0}
\right\}.
\]
For \(N>1\) on the polynomial branch, set
\[
\mathfrak U_{n,d,q}
=
\min\left\{
4,\,
b_{n,d,q}^2+v_{n,d,q}+4e^{-n\gamma_0}
\right\}.
\]
Then, with \(\tau(P)\), the average treatment effect, defined in \cref{obj:def:ate-functional}, and \(r_{n,d,q}\), the frontier rate, defined in \cref{obj:def:frontier-rate},
\[
E_P\!\left[
\left(
\widehat\tau^{\mathrm{mix}}_{n,d,q}(\mathbf O_n)-\tau(P)
\right)^2
\right]
\le
\mathfrak U_{n,d,q}
\le
\overline C\,r_{n,d,q}.
\]
The deterministic estimator also satisfies
\[
E_P\!\left[
\left(
\widehat\tau^{\mathrm{mix}}_{n,d,q}(\mathbf O_n)-\tau(P)
\right)^2
\right]
\le
E_{P,\omega}\!\left[
\left(
\widetilde\tau(\mathbf O_n;\omega)-\tau(P)
\right)^2
\right],
\]
where the auxiliary expectation uses the independent Poisson prefix and uniform stream assignments, including the zero output when the prefix exceeds \(n\), specified in \cref{obj:def:mixed-count-estimator}.
\end{theoremv}

\begin{algorithmv}{def:activation-lower-handle}[Activated fuzzy-hypothesis experiment]
The activated handle \(\mathfrak H_{\mathrm{act}}\) is the ordered pair of augmented \(n\)-record laws constructed below, with the following inputs:
\begin{itemize}
\item \textbf{(Parameters.)} A constant \(\eta>0\), fixed before integers \(n,d\ge1\) and an arrival floor \(0<q\le1\), satisfying \cref{obj:ass:rare-arrival-slice}. Define the effective sample size, logarithmic scale, moment-matching degree, support endpoint, nominal rare-cell mass, and rare-cell count by
\[
N=nq,\qquad
\ell=\log(e+N),\qquad
K=\lceil\ell\rceil,\qquad
H=K^2,\qquad
b_0=\frac{\eta}{N\ell},\qquad
J=\min\left\{d-1,\left\lfloor\frac{1}{2b_0}\right\rfloor\right\}.
\]
\item \textbf{(Finite priors.)} An ordered pair \((\Pi_0,\Pi_1)\) of probability laws, each assigning probability one to a finite subset of \([1,H]\).
\item \textbf{(Moment matching and separation.)} The priors satisfy
\[
\int z^v\,d\Pi_0(z)=\int z^v\,d\Pi_1(z)
\qquad\text{for every integer }0\le v\le K,
\]
and
\[
\left|\int z^{-1}\,d\Pi_1(z)-\int z^{-1}\,d\Pi_0(z)\right|
\ge\frac1{12}.
\]
\end{itemize}

For each \(i\in\{0,1\}\), independently draw a latent coordinate \(Z_x\sim\Pi_i\) for every baseline index \(x\in\{0,\ldots,d-1\}\). Conditional on this single latent vector, generate \(n\) independent records using the following common rule. Draw the baseline index with probabilities
\[
P(X=x)=
\begin{cases}
b_0,&0\le x<J,\\
1-Jb_0,&x=J,\\
0,&J<x<d.
\end{cases}
\]
The first \(J\) indices are the rare cells, and index \(J\) is the reservoir. Draw treatment independently with
\[
A\sim\operatorname{Bernoulli}(1/2),
\qquad S=0.
\]
Independently of the latent vector, baseline index, and treatment, draw and reveal the activation flag
\[
\mathsf B_{\mathrm{act}}\sim\operatorname{Bernoulli}(qH).
\]
For a rare cell \(x\), the same coordinate \(Z_x\) governs both treatment arms and all records in that cell.

If \(\mathsf B_{\mathrm{act}}=0\), set \((R,RY)=(0,0)\). Conditional on \(\mathsf B_{\mathrm{act}}=1\), draw \((R,RY)\) according to
\[
\begin{array}{c|ccc}
 & (R,RY)=(1,1) & (R,RY)=(1,0) & (R,RY)=(0,0)\\ \hline
 A=1,\ X=x<J & 1/H & (Z_x-1)/H & 1-Z_x/H\\
 A=0,\ X=x<J & 0 & Z_x/H & 1-Z_x/H\\
 X\ge J,\ \text{either arm} & 0 & 1 & 0
\end{array}
\]
Each augmented record is
\[
\bigl(\mathsf B_{\mathrm{act}},(X,A,S,R,RY)\bigr).
\]
The \(i\)-th component of \(\mathfrak H_{\mathrm{act}}\) is the law of these \(n\) augmented records after integrating over the product prior \(\Pi_i^{\otimes d}\). Thus the latent vector is drawn once for the entire sample, and the records are independent conditional on that vector. The handle is defined for every prior pair satisfying the stated conditions.
\end{algorithmv}

\begin{lemmav}{lem:reciprocal-best-approximation}[Reciprocal best uniform approximation]
For every integer \(K\ge 2\), set \(H=K^2\). Then
\[
\inf_{\substack{p\ \mathrm{real\ polynomial}\\ \deg(p)\le K}}
\sup_{z\in[1,H]}\left|z^{-1}-p(z)\right|
=
\frac{H-1}{2H}
\left(\frac{K-1}{K+1}\right)^K.
\]
\end{lemmav}

\begin{lemmav}{lem:moment-matched-reciprocal-priors}[Moment-matched reciprocal priors]
For every integer \(K\ge 2\), set \(H=K^2\). There exist finitely supported probability laws \(\Pi_0\) and \(\Pi_1\) on \([1,H]\) such that
\[
\int z^v\,d\Pi_0(z)=\int z^v\,d\Pi_1(z)
\qquad\text{for every integer }v\in\{0,\ldots,K\},
\]
and
\[
\left|\int z^{-1}\,d\Pi_1(z)-\int z^{-1}\,d\Pi_0(z)\right|
\ge \frac{1}{12}.
\]
\end{lemmav}

\begin{lemmav}{lem:one-cell-testing-family}[One-cell testing family bounds]
There exist a universal constant \(\underline c>0\) and constants \(\underline c_\alpha>0\) for every \(\alpha\in(0,1)\), independent of \(n,d,q\), such that the following holds whenever:
\begin{itemize}
\item \textbf{(Sample and alphabet sizes.)} \(n,d\) are integers with \(n\ge1\) and \(d\ge1\).
\item \textbf{(Arrival probability.)} \(0<q\le1\).
\item \textbf{(Rare-arrival slice.)} The restriction in \cref{obj:ass:rare-arrival-slice} holds.
\end{itemize}
Write \(N=nq\) for the effective sample size. There exists a family of full-data laws \(\{P_u:u\in[1/4,3/4]\}\), chosen independently of \(\alpha\), such that, for every \(u\in[1/4,3/4]\),
\[
P_u\in\mathcal M_{n,d,q},\qquad
P_u\!\left(X=0,\ S(0)=S(1)=0,\ Y(0)=0\right)=1,
\qquad
P_u(R=1)=q,\qquad
\tau(P_u)=u.
\]
Here \(X=0\) denotes the first baseline category in the alphabet \(\{0,\ldots,d-1\}\); the model class \(\mathcal M_{n,d,q}\) and average treatment effect \(\tau(P)\) are defined in \cref{obj:def:model-class,obj:def:ate-functional}.

There exist distinct \(u_0,u_1\in[1/4,3/4]\) such that every possibly randomized estimator \(\mathsf T\in\mathcal T_n\), as defined in \cref{obj:def:minimax-risk}, satisfies
\[
\underline c\min\{1,N^{-1}\}
\le
\max\left\{
\mathbb E_{P_{u_0}}\!\left[(\mathsf T-u_0)^2\right],
\mathbb E_{P_{u_1}}\!\left[(\mathsf T-u_1)^2\right]
\right\}.
\]
Moreover, the minimax squared-error risk in \cref{obj:def:minimax-risk} satisfies
\[
\mathfrak R_{n,d,q}\ge
\underline c\min\{1,N^{-1}\}.
\]

For every \(\alpha\in(0,1)\), there exist an integer \(M\ge1\) and pairwise distinct grid points \(u_0,\ldots,u_M\in[1/4,3/4]\) from this same family such that every uniformly honest connected-interval procedure \(\mathsf I\in\mathcal C_{n,d,q,\alpha}\), as defined in \cref{obj:def:interval-length-risk}, satisfies
\[
\underline c_\alpha\min\{1,N^{-1/2}\}
\le
\max_{i=0,\ldots,M}\mathbb E_{P_{u_i}}[u-l],
\]
where \((l,u)\) are the random endpoints returned by \(\mathsf I\). Moreover, the minimax expected interval length in \cref{obj:def:interval-length-risk} satisfies
\[
\mathfrak L_{n,d,q,\alpha}\ge
\underline c_\alpha\min\{1,N^{-1/2}\}.
\]
All expectations include the observed sample and the procedure's randomization.
\end{lemmav}

\begin{theoremv}{thm:full-experiment-lower}[Full-experiment minimax lower bound]
There exists a universal lower constant \(\underline c>0\) such that, for every \(n,d\in\mathbb N\) and \(q\in\mathbb R\) satisfying
\begin{itemize}
\item \textbf{(Sample size and dimension.)} \(n\ge1\) and \(d\ge1\).
\item \textbf{(Arrival parameter.)} \(0<q\le1\).
\item \textbf{(Rare-arrival slice.)} The restriction in \cref{obj:ass:rare-arrival-slice} holds.
\end{itemize}
the minimax squared-error risk \(\mathfrak R_{n,d,q}\) defined in \cref{obj:def:minimax-risk} obeys
\[
\mathfrak R_{n,d,q}\ge \underline c\,r_{n,d,q},
\]
where \(r_{n,d,q}\) is the frontier rate defined in \cref{obj:def:frontier-rate}.
\end{theoremv}
